\documentclass[10pt,letterpaper]{article}
\usepackage{cornell-notes}

\title{Clause 23.04.07: Suffix For Basic String Literals}
\author{}
\date{}
\setDocTitle{Clause 23.04.07: Suffix For Basic String Literals}
\setDocSubtitle{C++ 2024}
\setDocOwner{C++ 2024 Cornell Notes}
\setDocDate{\today}
\setCornellCollection{C++ 2024 Cornell Notes}
\setCornellUnitType{Clause}
\setCornellUnitNumber{23.04.07}
\setCornellUnitTitle{Suffix For Basic String Literals}
\hypersetup{pdftitle={Clause 23.04.07: Suffix For Basic String Literals},pdfsubject={Cornell notes},pdfauthor={}}

\begin{document}
\maketitle
\begin{CornellOverviewBox}{Overview}
\textbf{Classification:} Standard library - Normative\\[2pt]
\textbf{Learning objective:} Understand the rules, terminology, and semantic consequences associated with Suffix for basic\_string literals.
\end{CornellOverviewBox}\begin{CornellSummaryBox}{Summary}
{\sffamily\bfseries\color{Accent} Summary}\par\vspace{3pt}Section 23.4.7 focuses on Suffix for basic\_string literals. Apply the specified interface together with its mandates, effects, result conditions, exception behavior, and complexity constraints. When applying it, preserve the distinction between mandatory requirements, recommendations, permissions, and behavior deliberately left undefined, unspecified, or implementation-defined.
\end{CornellSummaryBox}

\end{document}
